\subsection{Loop-\texttt{else} Completion Semantics}

Python allows an \texttt{else} clause not only after \texttt{if}, but also after \texttt{for} and \texttt{while}. This surprises many programmers because the word \texttt{else} is usually associated with conditional branching. In loop syntax, however, \texttt{else} has a different meaning.

A loop-\texttt{else} block does not mean ``if the loop condition was false at the beginning.'' It means: run this block only if the loop finished normally.

For a \texttt{for} loop, normal completion means that the iterator was exhausted. For a \texttt{while} loop, normal completion means that the loop condition became false. If the loop is terminated by \texttt{break}, the \texttt{else} block is skipped.

\begin{quote}
A loop-\texttt{else} block runs after normal loop completion, not after early loop interruption.
\end{quote}

This makes loop-\texttt{else} useful for search patterns: if the loop finds the desired item, it uses \texttt{break}; if no item was found and the loop completed naturally, the \texttt{else} block reports the unsuccessful search.

\subsubsection{The Unique \texttt{else} Clause Semantics Applied to \texttt{while} and \texttt{for} Blocks}

The syntax is simple:

\begin{verbatim}
for item in items:
    loop_body
else:
    completion_body
\end{verbatim}

or:

\begin{verbatim}
while condition:
    loop_body
else:
    completion_body
\end{verbatim}

The \texttt{else} block is aligned with the loop header, not with an inner \texttt{if}. This indentation is essential. The \texttt{else} belongs to the loop.

A \texttt{for} example without \texttt{break}:

\begin{verbatim}
items = ["soup", "pretzels", "muffin"]

for item in items:
    print(f"checking: {item}")
else:
    print("loop completed normally")

print("after loop")
\end{verbatim}

Possible output:

\begin{verbatim}
checking: soup
checking: pretzels
checking: muffin
loop completed normally
after loop
\end{verbatim}

The iterator produced all values. No \texttt{break} interrupted the loop. Therefore the \texttt{else} block executed.

Now add a \texttt{break}:

\begin{verbatim}
items = ["soup", "pretzels", "muffin"]

for item in items:
    print(f"checking: {item}")

    if item == "pretzels":
        print("pretzels found, leaving loop")
        break
else:
    print("loop completed normally")

print("after loop")
\end{verbatim}

Possible output:

\begin{verbatim}
checking: soup
checking: pretzels
pretzels found, leaving loop
after loop
\end{verbatim}

The \texttt{else} block did not execute. The loop did not complete by exhausting the iterable. It was interrupted by \texttt{break}.

The control-flow shape is:

\begin{verbatim}
for loop starts
    -> get next item
    -> execute body
        if break occurs:
            jump after loop and skip else
    -> if iterator is exhausted:
            execute else
            continue after loop
\end{verbatim}

The same rule applies to \texttt{while}:

\begin{verbatim}
n = 1

while n <= 3:
    print(f"n: {n}")
    n = n + 1
else:
    print("while condition became false")

print("after loop")
\end{verbatim}

Possible output:

\begin{verbatim}
n: 1
n: 2
n: 3
while condition became false
after loop
\end{verbatim}

The loop stopped because \texttt{n <= 3} became false. This is normal completion, so the \texttt{else} block executed.

With \texttt{break}:

\begin{verbatim}
n = 1

while n <= 5:
    print(f"n: {n}")

    if n == 3:
        print("breaking at three")
        break

    n = n + 1
else:
    print("while condition became false")

print("after loop")
\end{verbatim}

Possible output:

\begin{verbatim}
n: 1
n: 2
n: 3
breaking at three
after loop
\end{verbatim}

The \texttt{else} block was skipped because the loop ended through \texttt{break}, not through the condition becoming false.

The loop-\texttt{else} rule is therefore not symmetric with \texttt{if}-\texttt{else}. In an \texttt{if} statement, \texttt{else} means ``the condition was false.'' In a loop, \texttt{else} means ``the loop was not interrupted by \texttt{break}.''

\begin{verbatim}
if condition:
    runs when condition is true
else:
    runs when condition is false

for item in items:
    runs for each produced item
else:
    runs when traversal finishes without break
\end{verbatim}

This is why the name \texttt{else} can feel misleading in loops. It is better to read it mentally as:

\begin{quote}
after normal loop completion
\end{quote}

\subsubsection{Conditional Execution: Triggering Logic Blocks Only After Non-Interrupted Loop Completion}

The most useful application of loop-\texttt{else} is search logic.

Without loop-\texttt{else}, programmers often use a flag:

\begin{verbatim}
names = ["Jerry", "Elaine", "Kramer"]
target = "Newman"
found = False

for name in names:
    print(f"checking: {name}")

    if name == target:
        found = True
        print(f"found: {name}")
        break

if not found:
    print("target was not found")
\end{verbatim}

Possible output:

\begin{verbatim}
checking: Jerry
checking: Elaine
checking: Kramer
target was not found
\end{verbatim}

The flag \texttt{found} records whether the loop ended because the target was found.

The loop-\texttt{else} version removes the flag:

\begin{verbatim}
names = ["Jerry", "Elaine", "Kramer"]
target = "Newman"

for name in names:
    print(f"checking: {name}")

    if name == target:
        print(f"found: {name}")
        break
else:
    print("target was not found")
\end{verbatim}

Possible output:

\begin{verbatim}
checking: Jerry
checking: Elaine
checking: Kramer
target was not found
\end{verbatim}

The meaning is direct: if the loop reaches a match, it breaks; if it never reaches a match, it completes normally and the \texttt{else} block runs.

When the target exists:

\begin{verbatim}
names = ["Jerry", "Elaine", "Kramer"]
target = "Elaine"

for name in names:
    print(f"checking: {name}")

    if name == target:
        print(f"found: {name}")
        break
else:
    print("target was not found")
\end{verbatim}

Possible output:

\begin{verbatim}
checking: Jerry
checking: Elaine
found: Elaine
\end{verbatim}

The \texttt{else} block is skipped because \texttt{break} interrupted the loop.

The list being searched is still an ordinary list object:

\begin{verbatim}
names = ["Jerry", "Elaine", "Kramer"]

print(f"type(names): {type(names)}")
print(f"dir(names) sample: {dir(names)[:10]}")
print(f"len(names): {len(names)}")
print(f"has __iter__: {'__iter__' in dir(names)}")
print(f"has __contains__: {'__contains__' in dir(names)}")
\end{verbatim}

Possible output:

\begin{verbatim}
type(names): <class 'list'>
dir(names) sample: ['__add__', '__class__', '__class_getitem__',
'__contains__', '__delattr__', '__delitem__', '__dir__', '__doc__',
'__eq__', '__format__']
len(names): 3
has __iter__: True
has __contains__: True
\end{verbatim}

The \texttt{\_\_iter\_\_()} method allows traversal. The \texttt{\_\_contains\_\_()} method supports membership testing with \texttt{in}. The loop-\texttt{else} structure does not change the list. It changes only how the program reacts to the traversal outcome.

The same pattern works with numeric search over a \texttt{range}:

\begin{verbatim}
target = 42

for n in range(1, 50):
    if n == target:
        print(f"found target: {n}")
        break
else:
    print("target was not reached")
\end{verbatim}

Possible output:

\begin{verbatim}
found target: 42
\end{verbatim}

If the target is outside the described range:

\begin{verbatim}
target = 100

for n in range(1, 50):
    if n == target:
        print(f"found target: {n}")
        break
else:
    print("target was not reached")
\end{verbatim}

Possible output:

\begin{verbatim}
target was not reached
\end{verbatim}

The range object completed traversal normally. No \texttt{break} occurred. Therefore the \texttt{else} block ran.

A \texttt{continue} statement does not suppress the loop-\texttt{else}. It interrupts only the current iteration, not the whole loop.

\begin{verbatim}
items = ["soup", "", "pretzels", ""]

for item in items:
    if item == "":
        print("empty item skipped")
        continue

    print(f"processed: {item}")
else:
    print("loop completed without break")
\end{verbatim}

Possible output:

\begin{verbatim}
processed: soup
empty item skipped
processed: pretzels
empty item skipped
loop completed without break
\end{verbatim}

The loop used \texttt{continue}, but it did not use \texttt{break}. Therefore the \texttt{else} block executed.

This gives an exact distinction:

\begin{itemize}
    \item \texttt{break} prevents loop-\texttt{else}.
    \item \texttt{continue} does not prevent loop-\texttt{else}.
    \item normal exhaustion of a \texttt{for} loop triggers loop-\texttt{else}.
    \item normal falsification of a \texttt{while} condition triggers loop-\texttt{else}.
\end{itemize}

An exception that escapes the loop also prevents normal loop completion:

\begin{verbatim}
items = ["10", "20", "bad", "30"]

try:
    for item in items:
        print(f"converting: {item}")
        number = int(item)
        print(f"number: {number}")
    else:
        print("all items converted")
except ValueError as err:
    print("conversion failed")
    print(f"type(err): {type(err)}")
    print(f"message: {err}")
\end{verbatim}

Possible output:

\begin{verbatim}
converting: 10
number: 10
converting: 20
number: 20
converting: bad
conversion failed
type(err): <class 'ValueError'>
message: invalid literal for int() with base 10: 'bad'
\end{verbatim}

The loop-\texttt{else} block did not run because the loop did not complete normally. The exception transferred control to the \texttt{except} block.

The practical reading is:

\begin{quote}
Use loop-\texttt{else} when the loop body contains a \texttt{break} that represents success, discovery, or early termination, and the \texttt{else} block represents the no-break case.
\end{quote}

A final compact pattern:

\begin{verbatim}
passwords = ["abc", "guest", "letmein"]
wanted = "swordfish"

for password in passwords:
    if password == wanted:
        print("access word found")
        break
else:
    print("Nobody expects the missing password.")
\end{verbatim}

Possible output:

\begin{verbatim}
Nobody expects the missing password.
\end{verbatim}

The loop inspected every candidate and never broke. The \texttt{else} block is therefore the completion route.

For control-flow graph analysis, loop-\texttt{else} adds one more outgoing edge from the loop structure: the normal-completion edge. A \texttt{break} follows the interrupted-exit edge and skips that completion block. This is the entire rule.